\section{The cyclic model and its coordinate meshes}\label{sec:setup}

We now choose the cyclic group and the two coordinate partitions on
which the two-color construction will be built.
All logarithms without an indicated base are natural. Throughout the cyclic
construction, $k$ is an integer exceeding an absolute threshold. Every
constant implicit in $O(\cdot)$ is absolute unless an explicit subscript
indicates otherwise, and every limiting assertion is as $k\to\infty$.
In particular, no threshold in the construction depends on the eventual
number of colors $r$.

Set
\begin{equation}\label{setup:parameters}
\begin{gathered}
 c=10^{-5},\qquad \delta=\frac1{10},\qquad \gamma=\frac1{100},\qquad
 D=\ceil{k^{1/10}},\qquad M=\ceil{k^{1/2}},\\
 h_0=D^2,\qquad p=k^{-1/20},\qquad H=k^{-2},\qquad
 \eta=\frac1{1000D}.
\end{gathered}
\end{equation}
We take $k$ large enough that $h_0<2M<k$.
The parameter $p$ will be the probability of a later perturbation. The
integer dilation
\begin{equation}\label{setup:dilation-definition}
 \lambda=\prod_{\substack{\ell\le h_0\\ \ell\text{ prime}}}
 \ell^{\floor{\log(2M)/\log\ell}}
\end{equation}
removes the small prime factors of all periods at most $2M$.
Least-common-multiple dilations also occur in the cyclic constructions
of Fox and Hunter~\cite{foxHunter2026}. The primary-factor choice here
is tailored to the growing dimension, as the next lemma records.

\begin{lemma}[Dilation of short periods]\label{setup:dilation}
We have $\log\lambda\le h_0\log(2M)$. If $1\le h\le2M$, then
every prime-power divisor of $h$ with prime at most $h_0$ divides
$\lambda$. Consequently, $h\mid\lambda$ whenever $h\le h_0$, and, for
every integer $t$, the denominator
\begin{equation}\label{setup:denominator}
 s=\frac{h}{\gcd(h,\lambda t)}
\end{equation}
is either $1$ or greater than $h_0$.
\end{lemma}

\begin{proof}
Each factor in \Cref{setup:dilation-definition} is at most $2M$,
and there are at most $h_0$ factors. If $\ell^a\mid h\le2M$, with
$\ell\le h_0$, then $a\le\floor{\log(2M)/\log\ell}$.
Thus the full $\ell$-primary part of $h$ divides $\lambda$, and hence
also $\gcd(h,\lambda t)$. All prime divisors of $s$ therefore exceed
$h_0$. This proves the assertions, including the case $t=0$, for which
$s=1$.
\end{proof}

\begin{lemma}[An elementary prime bound]\label{setup:prime}
For all sufficiently large integers $k$, there is a prime $P$ satisfying
$k<P\le2k^2$.
\end{lemma}

\begin{proof}
Put $n=k^2$. For a prime $\ell$, write $v_\ell(a)$ for the exponent
of $\ell$ in a positive integer $a$. Counting multiples of each power
of $\ell$ in the factorials gives
\[
 v_\ell\binom{2n}{n}
 =\sum_{a\ge1}\left(\floor{\frac{2n}{\ell^a}}
              -2\floor{\frac n{\ell^a}}\right)
 \le\floor{\log_\ell(2n)},
\]
because each summand is $0$ or $1$ and vanishes for $\ell^a>2n$.
If every prime factor of $\binom{2n}{n}$ were at most $k$, its value
would consequently be at most $(2n)^k$. On the other hand, the central
binomial coefficient is the largest of the $2n+1$ coefficients in
$(1+1)^{2n}$, so
\[
 \binom{2n}{n}\ge\frac{4^n}{2n+1}.
\]
For $n=k^2$, the logarithm of this lower bound is
$k^2\log4-\log(2k^2+1)$, whereas that of $(2n)^k$ is
$k\log(2k^2)$. The former is larger for all sufficiently large $k$.
There is therefore a prime divisor $P>k$; it is at most $2n$ because it
divides $(2n)!$.
\end{proof}

Choose such a prime $P$, and let $q$ be the least power of $P$ that is
at least $k^{ck/D}$. Minimality gives
\begin{equation}\label{setup:q-bounds}
 k^{ck/D}\le q\le2k^2 k^{ck/D},\qquad
 ck\log k\le D\log q\le ck\log k+D\log(2k^2).
\end{equation}
Here $q$ is a prime power, and need not itself be prime. Since
$h_0<k<P$ for large $k$, \Cref{setup:dilation-definition} shows that
$\gcd(\lambda,q)=1$. Define the remaining scales by
\begin{equation}\label{setup:mesh-parameters}
 H_x=\frac H\lambda,\qquad \alpha=q^{-(D+1)}.
\end{equation}
The parameter estimates used below include
\begin{equation}\label{setup:orders}
 D=O(k^{1/10}),\quad M=O(k^{1/2}),\quad
 \log\lambda=O(k^{1/5}\log k),\quad
 \log q=O(k^{9/10}\log k).
\end{equation}
These follow directly from \Cref{setup:parameters,setup:dilation,setup:q-bounds}.

Write $N=q^D$ and $G=\Z/N\Z$. We seek a two-coloring of $G$ with no
monochromatic progression
\begin{equation}\label{setup:cyclic}
 n_j=n_0+jd\quad(0\le j<k),\qquad n_0,d\in G,\qquad d\ne0.
\end{equation}
Every nonzero element of this prime-power cyclic group has order at
least $P>k$, so the $k$ displayed terms are distinct. On the circle
$\T=\R/\Z$, equipped with its usual distance and circumference one,
use the homomorphisms
\begin{equation}\label{setup:coordinates}
 x_i(n)=\frac n{q^i}\pmod1,\qquad
 y_i(n)=\lambda x_i(n)\qquad(1\le i\le D).
\end{equation}
They are well defined on $G$ since $q^i\mid N$. Write
$x(n)=(x_i(n))_{i=1}^D$ and $y(n)=(y_i(n))_{i=1}^D$, and fix their
representatives
\[
 \hat x(n)\in[0,1)^D,\qquad \tilde y(n)\in[-1/2,1/2)^D.
\]
The first coordinates will distinguish residue classes of a rational
return period. The dilation in the second coordinates removes rational
motion for periods at most $h_0$; the mesh near their representative
cut will then allow us to prove a global affine lift in that case.
The cut for the second representative system is the circle point
$1/2=-1/2$. Its distance function is
\begin{equation}\label{setup:cut-distance}
 \rho(Y)=\dist_{\T}(Y,1/2).
\end{equation}
Thus $0\le\rho\le1/2$, and $\rho$ is $1$-Lipschitz.

The first partition, denoted $U$, consists of the intervals
\[
 [aH_x,(a+1)H_x),\qquad 0\le a<k^2\lambda.
\]
This is an exact partition of $[0,1)$ because $H_x^{-1}=k^2\lambda$
is integral. The second partition becomes finer near its cut.
The positive scale $\alpha$ keeps this mesh finite; its choice will
be compared with the discrete step lattice in the short-period
argument of \Cref{dichotomy:main}. Put
\begin{equation}\label{mesh:definition}
 R=\log\left(1+\frac1{2\alpha}\right),\qquad
 m=\ceil{k^2R},\qquad
 a_\ell=\alpha\bigl(e^{\ell R/m}-1\bigr)\quad(0\le\ell\le m).
\end{equation}
The partition $V$ of $[-1/2,1/2)$ consists of
\begin{equation}\label{mesh:intervals}
 [-1/2+a_\ell,-1/2+a_{\ell+1}),\qquad
 [1/2-a_{\ell+1},1/2-a_\ell)
 \quad(0\le\ell<m).
\end{equation}
Since $a_0=0$ and $a_m=1/2$, these $2m$ intervals partition the
representative interval. They assign zero to the interval on its
right and the cut to the interval beginning at $-1/2$.
The same intervals also denote their images in $\T$. Closures and
neighborhoods of intervals will be taken on the circle; $|I|$ denotes
arc length.

For a circle point, $U(X)$ or $V(Y)$ denotes its individual interval
label, with the half-open conventions above. The full-label space and
label maps are
\begin{equation}\label{setup:labels}
 \begin{split}
 \cL&=U^D\times V^D,\\
 L(X,Y)&=\bigl((U(X_i))_{i=1}^D,(V(Y_i))_{i=1}^D\bigr),\\
 L(n)&=L(x(n),y(n)).
 \end{split}
\end{equation}
In particular, the identity of a label specifies each interval itself,
including its location on the circle.

\Needspace{3\baselineskip}
\begin{lemma}[Widths of the adaptive mesh]\label{mesh:width}
For every $I\in V$ and every $Y\in\overline I$,
\begin{equation}\label{mesh:width-bounds}
 \frac14H\bigl(\rho(Y)+\alpha\bigr)
 \le |I|\le2H\bigl(\rho(Y)+\alpha\bigr)\le2H.
\end{equation}
Moreover, $|V|=2m=O(k^3\log k)$.
\end{lemma}

\begin{proof}
For large $k$ we have $\alpha\le1/4$ and $R\ge1$. If
$\theta=R/m$, the bounds $k^2R\le m\le k^2R+1\le2k^2R$ give
$H/2\le\theta\le H$. On either interval indexed by $\ell$ in
\Cref{mesh:intervals}, the quantity $\rho(Y)+\alpha$ ranges, on the
closure, between $b=\alpha e^{\ell\theta}$ and $be^\theta$, while
$|I|=b(e^\theta-1)$. Hence
\[
 1-e^{-\theta}\le\frac{|I|}{\rho(Y)+\alpha}\le e^\theta-1.
\]
For $0\le\theta\le1$, the outside expressions are bounded below by
$\theta/2$ and above by $2\theta$, respectively. This proves the
first two inequalities in \Cref{mesh:width-bounds}; the last follows
from $\rho(Y)+\alpha\le3/4<1$.
Finally,
\[
 R\le(D+1)\log q=O(k\log k)
\]
by \Cref{setup:q-bounds}, so $2m\le2k^2R+2=O(k^3\log k)$.
\end{proof}

\begin{lemma}[Uniform neighborhoods]\label{mesh:neighbors}
For $I\in V$, let $w=|I|$ and
\[
 N(I)=\{Y\in\T:\dist_{\T}(Y,\overline I)\le w\}.
\]
Every $J\in V$ whose closure meets $N(I)$ satisfies $|J|\ge w/16$.
The arc $N(I)$ contains at most $49$ mesh endpoints.
\end{lemma}

\begin{proof}
Fix $Y_0\in\overline I$ and put $r_0=\rho(Y_0)+\alpha$.
Every $Y\in N(I)$ has circle distance at most $2w$ from $Y_0$.
By \Cref{mesh:width} and the Lipschitz property of $\rho$,
\[
 \rho(Y)+\alpha\ge r_0-2w\ge(1-4H)r_0\ge r_0/2,
\]
where the last inequality holds for $H\le1/8$. Apply the lower width
bound at a point of $\overline J\cap N(I)$ to obtain
\[
 |J|\ge\frac18Hr_0\ge\frac{w}{16}.
\]
The arc $N(I)$ has length $3w<1$. In its linear ordering, successive
mesh endpoints enclose a full mesh interval, of length at least
$w/16$ by the assertion just proved. If there are $b$ such endpoints,
then $(b-1)w/16\le3w$, so $b\le49$.
\end{proof}

For local counting in rotating coordinates, we ignore the fine mesh
within distance $\eta$ of the cut. Define the truncated output by
\begin{equation}\label{mesh:truncation-definition}
 V_\eta(Y)=
 \begin{cases}
 V(Y),&\rho(Y)\ge\eta,\\
 *,&\rho(Y)<\eta,
 \end{cases}
\end{equation}
where $*$ is a symbol distinct from every interval label. In particular,
equality at the threshold gives the actual interval label.

\begin{lemma}[Local complexity after truncation]\label{mesh:truncated}
On every circle arc of length at most $4H$, the output $V_\eta$ is
constant between at most $16004D$ breakpoints. Its value at a
breakpoint is also determined, and need not equal either neighboring
value.
\end{lemma}

\begin{proof}
The portion of the arc on which $\rho\ge\eta$ has at most two
components. Within one component, successive mesh endpoints have
spacing at least $H\eta/4$, by \Cref{mesh:width} applied at either
endpoint of the intervening mesh interval. The number of mesh
endpoints in these components is therefore at most
$4H/(H\eta/4)+2=16/\eta+2$. Add the at most two threshold points
$\rho=\eta$. Inside $\{\rho<\eta\}$ the output is constant, so no
other endpoints are needed. This gives at most
$16/\eta+4=16000D+4\le16004D$ breakpoints. At each breakpoint use
the half-open convention for $V$ and the weak inequality in
\Cref{mesh:truncation-definition}; this includes coincident mesh and
threshold points.
\end{proof}
